\part{Интегралы с переменными верхними пределами}

\section{Интегралы с переменными верхними пределами}
\begin{definition}
    пусть $f \in R[a,b] \forall x \in [a,b] \Rightarrow f \in R{[a,x]}$
    \[
    \int_a^x f(t) \, dt = F(x)
    \]
\end{definition}

\begin{paint}
    \begin{tikzpicture}[scale=1.2]

    % Оси
    \draw[->] (-0.5,0) -- (5,0) node[right] {$x$};
    \draw[->] (0,-0.5) -- (0,3.5) node[above] {$y$};

    % Точки a и b
    \coordinate (A) at (1,0);
    \coordinate (B) at (4,0);


    % График функции (кривая)
    \draw[thick, red]
        (1,1) 
        .. controls (1.5,2.2) and (2.5,2.3) .. (3,2)
        .. controls (3.5,1.8) and (3.8,2.2) .. (4,2.4);

    % Заливка области
    \fill[blue!20]
        (1,0) -- (1,1)
        .. controls (1.5,2.2) and (2.5,2.3) .. (3,2)
        .. controls (3.5,1.8) and (3.8,2.2) .. (4,2.4)
        -- (4,0) -- cycle;
        
    \fill[pattern=north east lines]
        (2.5,0) -- (2.5,2.16)
        .. controls (3,2) and (3.5,1.8) .. (4,2.4)
        -- (4,0) -- cycle;

    % Надпись f(x)
    \node at (2.5,3) {$f(x)$};

    % Вертикальные линии
    \draw[thick] (1,0) -- (1,1);
    \draw[thick] (4,0) -- (4,2.4);
    \draw[thick] (2.5,0) -- (2.5,2.16);

    % Точка O
    \node at (-0.2,-0.2) {$O$};
    \node at (1,-0.2) {$a$};
    \node at (4,-0.2) {$b$};
    \node at (2.5,-0.2) {$c$};

    \end{tikzpicture}
\end{paint}

\begin{theorem}
    $ f \in R[a,b] \Rightarrow \in \mathbb{C}[a,b] $
    \begin{proof}
        $ \forall  x \in [a,b]$ x - фиксированное число
        $x + n \in[a,b] $, где n - приращение
        $ \delta F = F(x + n) - F(x) = \int_a^{x+n} f(t) \, dt - \int_a^x f(t) \, dt = \int_x^{x+n} f(t) \, dt $ \\
        $ f \in R[a.b] f$ - ограниченная функция, значит $ \exists M \geqslant 0: |f(x)| \leqslant M, \forall x \in [a,b] $ \\
        $ 0 \leqslant |\delta F(x)| = |\int_x^{x+n} f(t) \, dt| \leqslant k * |n| \space (|n| \to 0) \Rightarrow \lim\limits_{n \to 0} \delta F(x) = 0$

    \end{proof}
\end{theorem}